\section{Lesson 7: Reading Real Metaphysics}

In the earlier lessons you practiced the six questions one by one. Now all six work together on full pages of serious metaphysics across four stages of increasing independence: from a complete worked reading on the first passage, through guided reconstruction and rival comparison beside the One-Page Analysis Sheet (Appendix B), to an unaided reading on the fourth. These six headings remain a modern reading habit, never Aristotle's or Aquinas's own outline; once the questions become second nature, the headings fall away.

\subsection{Worked Passage 1: Causal Regress in Aquinas (\emph{ST} I, q.\,46, a.\,2, ad 7)}

Begin with a full six-question reading in canonical order---followed by an objection at strength, an honest partial concession, and a one-step defective near-neighbour---before taking over more of the work yourself on the next three passages.

In \emph{Summa Theologiae} I, q.\,46, a.\,2 (\emph{Utrum mundum incepisse sit articulus fidei}---``Whether it is an article of faith that the world began''), St.\ Thomas Aquinas holds that the world's temporal beginning is known by faith alone and cannot be demonstrated by natural reason. Objection 7 counters that if the world were eternal, human generation would run backward to infinity; since a father is his son's efficient cause (Aristotle, \emph{Physics} II), there would be an infinite regress of efficient causes, disproved in \emph{Metaphysics} II. Here is Aquinas's reply in the Leonine text (\emph{ST} I, q.\,46, a.\,2, ad 7, Rome, 1888), unmarked:

\begin{quote}
Ad septimum dicendum quod in causis efficientibus impossibile est procedere in infinitum per se; ut puta si causae quae per se requiruntur ad aliquem effectum, multiplicarentur in infinitum; sicut si lapis moveretur a baculo, et baculus a manu, et hoc in infinitum. Sed per accidens in infinitum procedere in causis agentibus non reputatur impossibile; ut puta si omnes causae quae in infinitum multiplicantur, non teneant ordinem nisi unius causae, sed earum multiplicatio sit per accidens; sicut artifex agit multis martellis per accidens, quia unus post unum frangitur. Accidit ergo huic martello, quod agat post actionem alterius martelli. Et similiter accidit huic homini, inquantum generat, quod sit generatus ab alio, generat enim inquantum homo, et non inquantum est filius alterius hominis; omnes enim homines generantes habent gradum unum in causis efficientibus, scilicet gradum particularis generantis. Unde non est impossibile quod homo generetur ab homine in infinitum. Esset autem impossibile, si generatio huius hominis dependeret ab hoc homine, et a corpore elementari, et a sole, et sic in infinitum.
\end{quote}

A close English rendering, matching the Latin clause by clause:

\begin{quote}
To the seventh it must be said that in efficient causes it is impossible to proceed to infinity per se; for instance, if the causes which are required per se for a certain effect were multiplied to infinity, as if a stone were moved by a stick, and the stick by a hand, and this to infinity. But per accidens to proceed to infinity in acting causes is not judged impossible; for instance, if all the causes which are multiplied to infinity hold the order of only one cause, but their multiplication is per accidens, as an artisan acts with many hammers per accidens because one breaks after another. It happens incidentally, therefore, to this hammer that it acts after the action of another hammer. And similarly it happens incidentally to this man, insofar as he begets, that he was begotten by another, for he begets insofar as he is a man, and not insofar as he is the son of another man; for all men who beget hold one grade in efficient causes, namely the grade of a particular generator. Hence it is not impossible that a man be begotten by a man to infinity. It would, however, be impossible if the generation of this man depended on this man, and on an elementary body, and on the sun, and so on to infinity.
\end{quote}

\subsubsection*{Claim}

In efficient causes, an infinite regress \emph{per se} is impossible, whereas an infinite regress \emph{per accidens} is not judged impossible; hence an endless line of fathers begetting sons is not ruled out by the ban on infinite regress.

\emph{Plain-language reading.} When an effect needs a chain of causes each acting only because the next drives it---stone moved by stick, stick by hand---that chain cannot run backward forever. Nothing forbids an endless line of causes filling one single role in succession, like a craftsman replacing broken hammers: whether an earlier hammer snapped makes no difference to how this hammer drives a nail. Likewise, when a man fathers a child, his own having been fathered makes no difference to his act of begetting; he begets because he is a man, not because he is a son. All human fathers occupy one causal tier, so an endless human lineage is not barred here. What \emph{would} be impossible is an endless upward stack of distinct tiers all required for this birth: this father, an elementary body, the sun, and so on without end.

We kept \emph{per se} and \emph{per accidens} in Latin because flattening them into ``simultaneous'' and ``accidental'' distorts the claim: being begotten by one's father is no mishap (it is incidental strictly \emph{to his present act of begetting}), and simultaneous events need not derive causal power from one another.

\subsubsection*{Meaning}

Four terms must stay fixed. First, \emph{per se} and \emph{per accidens} of causes (not predication from the \emph{Posterior Analytics}): causes are ordered \emph{per se} when each subordinate cause acts only by actively borrowing power from the higher cause (stick and hand); they are ordered \emph{per accidens} when earlier causes do not enter into what the present cause needs to act (breaking hammers). Broaden \emph{per se} to any historical prerequisite, and every grandfather becomes a \emph{per se} cause, restoring Objection 7; read \emph{per accidens} as ``random,'' and human generation fails immediately. Second, \emph{holding the order of only one cause} (\emph{non teneant ordinem nisi unius causae}): distinct individuals filling one functional slot. Read it as numerical singularity and the hammers fail; dilute it to any chain with one outcome and hand-stick-stone merges into it. Third, \emph{grade} (\emph{gradus}): a distinct rank of causal power in a hierarchy (particular generator, elementary body, sun). Take \emph{grade} as genealogical distance and every generation becomes a new tier. Because the closing ladder reflects medieval astronomy, keeping \emph{grade} alive requires showing that real hierarchies of concurrent dependence exist in nature. Fourth, \emph{inquantum} (``insofar as''): isolates the formal respect under which an agent acts (\emph{inquantum homo}, not \emph{inquantum filius}). Drop \emph{inquantum}, and Aquinas denies the truism that no man exists without a father.

\subsubsection*{Warrant}

Provenance and revelation remain strictly apart from natural warrant. \textbf{This citation establishes only that St.\ Thomas Aquinas held this distinction and offered this reply; it supplies no reason at all to believe that an endless \emph{per se} series is impossible or that real causal series divide this way.} In the body of q.\,46, a.\,2, Aquinas labels the world's temporal beginning as \textbf{revelation} (\emph{sola fide tenetur}, Genesis 1:1 and the Creed): we neither launder faith into philosophy nor demand natural proof of a revealed article. Inside Reply 7, Aquinas argues by natural reason alone.

Three moves braid the paragraph's warrant. First, \textbf{experience} paired with \textbf{definition} anchors the contrast: observation shows a stick stopping when the hand stops, whereas a fresh hammer strikes after an old one breaks and a son fathers children after his father dies; defeating Objection 7 then follows as a \textbf{logical consequence}. Second, for \emph{impossibile est procedere in infinitum per se}, Reply 7 gives no proof on this page; it draws on a \textbf{prior result} argued in \emph{Summa contra Gentiles} I, c.\,13, nn.\,14, 15, and 33 (ordered movers halting without a first in n.\,14; instrumental movers lacking a principal mover in n.\,15; ordered efficient causes collapsing without a first in n.\,33; compressed in \emph{ST} I, q.\,2, a.\,3). Third, until the reader checks \emph{SCG} I, c.\,13, that rule stands here on \textbf{provenance only}---a diagnosis, not a warrant. Finally, the closing conditional (man, elementary body, sun) invokes a \textbf{metaphysical principle} of subordinate natural agency wedded to medieval cosmology, unargued here.

\subsubsection*{Burden}

Reply 7 performs two jobs: a narrow dialectical job inside Question 46 (refuting Objection 7 by placing human fathers on one \emph{per accidens} tier) and a wider metaphysical job whenever this rule supports cosmological proofs. Test two candidates openly with the hostile-reader question: \emph{could an intelligent opponent accept every definition and illustration here and still deny this proposition?}

\emph{First candidate (rejected):} ``In efficient causes it is impossible to proceed to infinity \emph{per se}.'' An opponent can grant this completely for chains of wholly passive sticks---bordering on unpacking ``wholly derived''---and yet concede nothing about nature or a first cause unless some real natural series actually works like hand, stick, and stone.

\emph{Second candidate (the burden-bearing premise):} Pointed to conditionally in the closing sentence (\emph{esset autem impossibile, si\dots}), the metaphysical bridge is that \textbf{some real series of causal dependence in nature is ordered \emph{per se} in this derived-causing sense---so that borrowed agency governs actual tiers of natural dependence beyond local mechanical tools.} An opponent can accept every definition, agree about hammers, fathers, and sticks, and still deny without contradiction that particles, organisms, or substances depend vertically on concurrent higher causes. Aquinas writes the closing sentence as an \emph{if}-clause and does not prove this bridge here; isolating it locates the real work before cosmological arguments begin.

\subsubsection*{Consequence}

Granting only what Reply 7 itself asserts defeats Objection 7 and establishes a conditional rule: \emph{if} a real series of efficient causes is ordered \emph{per se} in the derived-causing sense, it cannot run endlessly without an underived member in that series. Granting \emph{also} the burden-bearing premise for a specific natural series yields that \emph{that} series ends in a first underived member---and nothing stronger.

Here, and nowhere else in Lesson 7, two displayed lines earn their place because ordinary English hides both a relation ambiguity and a quantifier-order shift inside ``every cause depends on another.'' Let quantifiers range over members of a single causal series, and write \linline{\logicop{Dep}(x, y)} for \emph{the causing done by \linline{x} is derived from the causing done by \linline{y}, so that \linline{x} causes only because \linline{y} causes} (reading (2) over all subordinate dependent members, since a first cause does not derive its causing from itself). Compare:
\ldisplay{(1) \lquad \lall x\lthin \lex y\lthin \logicop{Dep}(x, y)}
\ldisplay{(2) \lquad \lex y\lthin \lall x\lthin \logicop{Dep}(x, y)}
These lines expose two features at a glance. First, \textbf{ambiguity in the relation}: read \linline{\logicop{Dep}(x, y)} loosely as historical succession (``\linline{x} was caused by \linline{y}''), and (1) is what the \emph{per accidens} concession allows without end; read it strictly as derived causing, and (1) is what the \emph{per se} rule forbids. Second, \textbf{quantifier order}: (1) says each \linline{x} depends on some \linline{y}, whereas (2) says some \emph{one} head \linline{y} sustains every dependent member. Cosmological proofs aim at something like (2), yet Reply 7 asserts only that (1) under derived causing is impossible---without proving that any cosmos-wide series is governed by \linline{\logicop{Dep}(x, y)} strictly or that separate \emph{per se} series converge on one unique head.

\emph{Not yet established.} Granting everything Aquinas asserts in Reply 7 leaves unestablished: (1) a first cause of being, still less one God, creator of heaven and earth (even granting the bridge for a natural series yields only an underived member of that series, not a unique, simple, intelligent God); (2) that any natural chain beyond local mechanical tools is \emph{per se} ordered; and (3) whether the world began in time or existed eternally.

\subsubsection*{Denial}

Rejecting the burden-bearing premise is \textbf{rival framework}, and specifically \emph{not} \textbf{contradiction}. Denying that finite substances derive their present causal power vertically from higher causes breaks no law of logic; claiming ``incoherence'' hurts credibility where genuine contradiction obtains. The serious denier holds a \textbf{rival framework} in which natural systems act natively through intrinsic powers (mass, charge, biological organization) once brought into being, making cosmic chains resemble generation rather than sticks in hands. Simply declining the \emph{per se} / \emph{per accidens} vocabulary is \textbf{non-assent}. A Thomist may still urge an \textbf{explanatory cost}---asking what sustains contingent substances and their lawful powers---which the rival answers inside his framework or accepts as brute.

\subsubsection*{Objection at Full Strength}

Here is a constructed dilemma attacking the burden-bearing premise at full strength:

\begin{quote}
Either \emph{per se} ordering is defined narrowly enough to fit hand, stick, and stone, or broadly enough to govern natural substances generally. Narrowly defined, a stick is a passive tool pushed by contact at that instant; an endless stack of sticks cannot initiate motion, yet the only \emph{per se} series we encounter are local mechanical assemblies, while hot iron, adult fathers, and particles act by intrinsic properties with no visible hand behind them---leaving the rule true but cosmically sterile. Broadly defined to cover non-mechanical sustaining dependence, the concept gains cosmic reach, yet the proof against an endless descending hierarchy of sustaining ontological levels vanishes unless you pack ``borrowing from a single unborrowed source'' into the definition. Narrowly read, no cosmic series fits it; broadly read, its impossibility is stipulated rather than proved.
\end{quote}

\subsubsection*{Reply and Partial Concession}

The honest answer is a \textbf{partial concession}. Against \emph{ST} I, q.\,46, a.\,2, ad 7 alone, the dilemma lands: Reply 7 proves no universal \emph{per se} series of being, and rescuing it by defining ``finite being'' as ``instrument of a first cause'' would collapse into tautology. Yet inside Question 46, Aquinas needs only the negative half of his distinction to defeat Objection 7 on the narrow horn alone. Beyond Question 46, the dilemma names the exact address of the missing proofs: \emph{De ente et essentia} c.\,4, \emph{ST} I, q.\,2, a.\,3, and \emph{SCG} I, c.\,13. Reply 7 states the conditional rule, leaving both its proof for non-mechanical hierarchies and the proof of its antecedent to be earned elsewhere.

\keeppair
\subsubsection*{A Defective Near-Neighbour}

Read this fluent handbook exposition first without comment:

\begin{quote}
In any causal sequence we observe in nature, each efficient cause is itself caused by another: the flame is lit by the match, the match is struck by the hand, and the son is begotten by the father. Nothing can be the efficient cause of itself, since it would then have to exist prior to itself, which is absurd. Yet an infinite chain of efficient causes is impossible, because in any chain of causes, removing the first cause removes every intermediate cause and the final effect as well. Since every observed efficient cause in a given chain depends on another cause, that chain cannot stretch backward forever: it must terminate in a first efficient cause that is not itself caused by another.
\end{quote}

This paragraph differs from Aquinas's analysis by \textbf{one exact, locatable step}: it treats ``each efficient cause is caused by another''---line (1) read loosely as bare succession---as though it licensed the elimination rule, omitting the derived-causing premise separating \emph{per se} series from \emph{per accidens} succession. Blending hand-match with father-son and applying the elimination rule to ``any chain of causes'' commits the very mistake Reply 7 refutes. (Even once repaired for one chain, uniqueness across all chains remains a separate argument.)

\subsection{Worked Passage 2: Analogy and Divine Names in Aquinas (\emph{ST} I, q.\,13, a.\,3 and a.\,5)}

Read the two unmarked excerpts and the explicit prose reconstruction below, then draft your own answers to all six questions on the One-Page Analysis Sheet (Appendix B) before checking the analysis that follows.

In \emph{ST} I, q.\,13, a.\,3 (``Whether Any Name Can Be Applied to God in Its Literal Sense?'', 1920 English Dominican translation), Aquinas distinguishes what perfection-names signify from how they signify:

\begin{quote}
According to the preceding article, our knowledge of God is derived from the perfections which flow from Him to creatures, which perfections are in God in a more eminent way than in creatures. Now our intellect apprehends them as they are in creatures, and as it apprehends them it signifies them by names. \dots\ As regards what is signified by these names, they belong properly to God, and more properly than they belong to creatures, and are applied primarily to Him. But as regards their mode of signification, they do not properly and strictly apply to God; for their mode of signification applies to creatures.
\end{quote}

In Article 5 (``Whether What Is Said of God and of Creatures Is Univocally Predicated of Them?'', same translation), after explaining that inadequate effects receive in divided form what pre-exists unitedly in God (as the sun's single power produces varied forms, and as ``wise'' in man names a distinct quality while in God it leaves His essence uncomprehended), he concludes:

\begin{quote}
Univocal predication is impossible between God and creatures. \dots\ Neither, on the other hand, are names applied to God and creatures in a purely equivocal sense, as some have said. Because if that were so, it follows that from creatures nothing could be known or demonstrated about God at all; for the reasoning would always be exposed to the fallacy of equivocation. Such a view is against the philosophers, who proved many things about God, and also against what the Apostle says: ``The invisible things of God are clearly seen being understood by the things that are made'' (Rom.\ 1:20). Therefore it must be said that these names are said of God and creatures in an analogous sense, i.e.\ according to proportion. \dots\ Thus whatever is said of God and creatures, is said according to the relation of a creature to God as its principle and cause, wherein all perfections of things pre-exist excellently. Now this mode of community of idea is a mean between pure equivocation and simple univocation.
\end{quote}

Reconstructed into explicit prose first, the compressed argument runs on two rails. Against univocity: perfections united in God's simple essence appear divided in creatures; because we name perfections as grasped in creatures, the perfection signified (\emph{res significata}) belongs primarily to God while the creaturely mode of signifying (\emph{modus significandi}) does not, ruling out one identical defining account. Against pure equivocation: wholly unrelated senses would ruin every inference from creatures to God by equivocation, contradicting both the philosophers' proofs and Romans 1:20. Therefore perfection-names are predicated analogically (\emph{secundum analogiam, idest proportionem}), by the creature's causal relation to God, midway between pure equivocation and simple univocation (\emph{inter puram aequivocationem et simplicem univocationem}).

Stop here and fill out the One-Page Analysis Sheet (Appendix B) before comparing against the analysis below.

\subsubsection*{Claim}

Pure perfection-names apply properly to God in what they signify (\emph{res significata}) though not in their creaturely mode of signifying (\emph{modus significandi}), and are predicated of God and creatures neither univocally nor purely equivocally, but analogically by causal proportion. \emph{Plain-language reading.} Words like ``wise'' and ``good'' name realities present in God even more fully than in us, though framed in limited creaturely grammar; they neither share one exact definition across God and creatures nor split into unrelated homonyms, but link effects back to their divine source by proportion.

\subsubsection*{Meaning}

Fix \emph{res significata} (the perfection named) versus \emph{modus significandi} (its limited creaturely framing), \emph{univocal} (one identical defining account), \emph{purely equivocal} (wholly unrelated accounts), and \emph{analogical by one-to-another proportion} (as medicine is ``healthy'' by causing an animal's health, rather than two things pointing to a third standard above God). Blur analogy into metaphor and literal predication collapses; drop \emph{modus significandi} and divine simplicity dissolves.

\subsubsection*{Warrant}

Provenance and revelation sit beside natural premises here. Romans 1:20 enters as \textbf{revelation}---legitimate inside theology, yet kept strictly labelled and never converted into philosophical proof. Beside it, ``against the philosophers, who proved many things about God'' stands on \textbf{provenance only}---a diagnosis noting that citing past philosophers does not prove on this page that their proofs avoided equivocation. The philosophical steps rely on \textbf{prior results} (simplicity and eminence, qq.\,3--4 and q.\,13, a.\,4; demonstrability from creatures, q.\,2, aa.\,2--3, or a \textbf{hypothesis} until q.\,2 is checked), a \textbf{metaphysical principle} (inadequate effects receive divided likenesses, illustrated from \textbf{experience} via the medieval sun analogy), \textbf{logical consequence} (pure equivocation triggers the fallacy of equivocation), and \textbf{definition} (classifying proportionate naming).

\subsubsection*{Burden}

Once divine simplicity (q.\,3) and univocity are granted, excluding univocity is nearly automatic. The real burden-bearing premise of Article 5 is that \textbf{real causal proportion between creatures and God sustains a genuine middle semantic tier (\emph{secundum analogiam}) between univocity and pure equivocation, preserving valid inferential force without a shared univocal account.} An opponent can accept every definition here, grant simplicity, and still deny that any non-univocal middle avoids equivocation in a syllogism.

\subsubsection*{Consequence}

Granted simplicity, causal eminence, and the adequacy of causal analogy, pure perfection-names apply literally to God (as to \emph{res significata}) without pulling God into a creaturely genus or collapsing speech into homonyms. No notation appears here because the issue is semantic proportion, not scope.

Yet this doctrine creates a celebrated logical difficulty: because an analogical middle term does not carry one identical defining account (\emph{ratio}) across both premises of a syllogism, it threatens a \textbf{four-term fallacy}. Historically, this syllogistic problem is a construction of the later tradition---above all Thomas de Vio Cajetan in \emph{De nominum analogia} (1498)---not Aquinas's own theme in Question 13, which targets \emph{pure} equivocation rather than formalizing analogical syllogisms.

\emph{Not yet established.} Articles 3 and 5 do not establish: (1) how an analogical middle preserves strict syllogistic validity without a four-term fallacy or a thin univocal core; (2) positive insight into \emph{how} wisdom exists in God's simple essence beyond affirming the perfection and negating creaturely limits; or (3) that any specific perfection exists in God (presupposed from qq.\,2--12).

\subsubsection*{Denial}

\textbf{Rival framework}, not \textbf{contradiction}: either a univocity framework (Bl.\ John Duns Scotus below) insisting that valid demonstration demands at least one minimal univocal concept like \emph{being}, or a strict apophatic/agnostic framework restricting human concepts to finite creatures.

\subsection{Worked Passage 3: Comparing Burdens Across Rival Traditions (Scotus and Hume)}

Now read two unmarked passages outside the Thomistic line without prior reconstruction. Fill out the six headings of Appendix B on blank paper for each text before reading the comparison below.

The first passage is Bl.\ John Duns Scotus, \emph{Ordinatio} I, d.\,3, p.\,1, q.\,1--2, n.\,26 (Vatican edition):

\begin{quote}
Et ne fiat contentio de nomine univocationis, univocum conceptum dico, qui ita est unus, quod eius unitas sufficit ad contradictionem, affirmando et negando ipsum de eodem; sufficit etiam pro medio syllogistico, ut extrema unita in medio sic uno sine fallacia aequivocationis concludantur inter se uniri.
\end{quote}

A close English rendering of the Latin:

\begin{quote}
And lest there be a dispute over the word ``univocation,'' I call a concept univocal which is so unified that its unity suffices for contradiction when affirming and denying it of the same subject; it also suffices as a syllogistic middle term, so that extremes united in a middle term thus unified may be concluded to be united with one another without the fallacy of equivocation.
\end{quote}

The second is David Hume, \emph{Dialogues Concerning Natural Religion}, Part IX (1779, public domain), where Cleanthes replies to Demea's demand for an external cause of an eternal causal chain:

\begin{quote}
In such a chain, too, or succession of objects, each part is caused by that which preceded it, and causes that which succeeds it. Where then is the difficulty? But the whole, you say, wants a cause. I answer, that the uniting of these parts into a whole, like the uniting of several distinct countries into one kingdom, or several distinct members into one body, is performed merely by an arbitrary act of the mind, and has no influence on the nature of things. Did I show you the particular causes of each individual in a collection of twenty particles of matter, I should think it very unreasonable, should you afterwards ask me, what was the cause of the whole twenty. That is sufficiently explained in explaining the cause of the parts.
\end{quote}

Stop here and draft your own six-question analyses for both passages before reading on.

\subsubsection*{Scotus Compared with Aquinas (\emph{ST} I, q.\,13, a.\,5)}

\textbf{Claim and plain-language reading.} Scotus fixes a logical test for a univocal concept: unity sufficient for contradiction when affirmed and denied of the same subject, and sufficient to act as a syllogistic middle term (\emph{medium syllogisticum}) without equivocation. \textbf{Meaning.} Where Aquinas took univocity as one full generic defining account, Scotus narrows it to a minimal logical unity sufficient for contradiction and valid inference. \textbf{Warrant.} Within n.\,26 itself, the warrant is pure \textbf{definition} (\emph{univocum conceptum dico\dots}), openly stipulated to prevent verbal disputes.

\textbf{Burden and comparison.} By itself, n.\,26 carries no ontological burden---a clean conceptual standard. Scotus's substantive burden enters when he applies it to natural theology: \textbf{that no concept lacking this minimal univocal unity can carry a valid demonstration from creatures to God without a four-term fallacy, and that our intellect actually possesses such a minimally univocal concept (like \emph{being}) common to God and creatures without compromising divine simplicity.} Comparing burdens exposes the real clash between two Catholic doctors: Aquinas's burden (\emph{ST} I, q.\,13, a.\,5) is showing how causal analogy avoids a four-term fallacy without a univocal core; Scotus's burden (\emph{Ordinatio} I, d.\,3) is showing how minimal logical univocity avoids forcing God and creatures into one shared genus. \textbf{Consequence.} Granting n.\,26 gives a logical test; granting the broader premise requires at least one minimally univocal concept in natural theology. \emph{Not yet established} by n.\,26: that we actually form such a concept across God and creatures or that it leaves simplicity untouched. \textbf{Denial.} \textbf{Rival framework} (Aquinas's causal analogy), not contradiction.

\subsubsection*{Hume Compared with Aquinas (\emph{ST} I, q.\,46, a.\,2, ad 7)}

\textbf{Claim and plain-language reading.} In a succession where each object is caused by its predecessor, explaining every individual member sufficiently explains the whole; asking for an extra cause of ``the whole'' mistakes a mental grouping (like twenty particles of matter) for an extra reality. \textbf{Meaning.} Everything turns on \emph{succession of objects} (self-standing particulars in time) and \emph{whole} (``an arbitrary act of the mind'' with no effect on reality). If a causal whole is instead a real structure of concurrent dependence, the twenty-particles analogy breaks. \textbf{Warrant.} \textbf{Experience} and analogy (countries united into a kingdom, twenty particles) support a \textbf{metaphysical principle}---a collection is nothing over and above its members, so explaining every member leaves nothing unexplained---yielding Cleanthes's refusal of Demea's demand by \textbf{logical consequence}.

\textbf{Burden and comparison.} An opponent can grant that mentally grouping twenty marbles adds no twenty-first marble and still refuse Cleanthes's conclusion. The burden-bearing premise is that \textbf{the causal series of contingent beings behaves like a mental aggregation of twenty self-standing particles in temporal succession, leaving nothing about the collection's existence or concurrent activity unexplained once each predecessor is named.} Comparing burdens with \emph{ST} I, q.\,46, a.\,2, ad 7 shows both convergence and collision. On a \emph{per accidens} temporal succession (fathers or hammers), Aquinas and Hume partly agree that bare aggregation does not refute an endless chain, and Hume cleanly exposes the member-to-whole fallacy for mere aggregates. Their burdens collide squarely at \emph{per se} concurrent dependence: Hume assumes all cosmic causation works like twenty separate particles handed along in time, whereas Aquinas's cosmological arguments hold that finite beings stand in concurrent vertical dependence where no member owns its act of being natively. \textbf{Consequence.} Granted Cleanthes's premise, Demea's member-to-whole step fails for aggregates. \emph{Not yet established:} (1) that no genuine question remains why \emph{any} succession of contingent particles exists at all rather than nothing; or (2) that concurrent \emph{per se} dependence reduces to a mental grouping of twenty successive particles. \textbf{Denial.} \textbf{Rival framework} joined to a clash over \textbf{explanatory cost} (not contradiction): the Thomist presses the cost of leaving an entire contingent chain brute, while the Humean accepts that stopping point.

\subsection{Worked Passage 4: Unaided Exit Reading --- Aquinas on Creation by Participation (\emph{ST} I, q.\,44, a.\,1, corp.)}

Take a pencil and the One-Page Analysis Sheet (Appendix B), read the unannotated text below, and complete all six headings unaided before checking your work against the discussion that follows.

Here is the body of St.\ Thomas Aquinas, \emph{ST} I, q.\,44, a.\,1, ``Whether It Is Necessary That Every Being Be Created by God?'' (1920 English Dominican translation):

\begin{quote}
It must be said that every being in any way existing is from God. For whatever is found in anything by participation, must be caused in it by that to which it belongs essentially, as iron becomes ignited by fire. Now it has been shown above (Q.\,3, A.\,4) when treating of the divine simplicity that God is the essentially self-subsisting Being; and also it was shown (Q.\,11, AA.\,3, 4) that subsisting being must be one; as, if whiteness were self-subsisting, it would be one, since whiteness is multiplied by its recipients. Therefore all beings apart from God are not their own being, but are beings by participation. Therefore it must be that all things which are diversified by the diverse participation of being, so as to be more or less perfect, are caused by one First Being, Who possesses being most perfectly.

Hence Plato said (Parmen.\ xxvi) that unity must come before multitude; and Aristotle said (Metaph.\ ii, text 4) that whatever is greatest in being and greatest in truth, is the cause of every being and of every truth; just as whatever is the greatest in heat is the cause of all heat.
\end{quote}

Stop here. Write out your six headings on blank paper before reading the solution below.

\subsubsection*{Claim}

Every being existing in any way is caused by God, because whatever possesses existence by participation rather than essentially must be caused by the one self-subsisting Being. \emph{Plain-language reading.} Nothing existing stands outside God's causal production: just as heated iron receives heat from fire because it is not fire itself, every finite thing shares in existence without being pure existence itself, and must therefore be caused by the one Being whose essence is existence.

\subsubsection*{Meaning}

Fix \emph{by participation} versus \emph{essentially} (receiving a limited share of a perfection distinct from one's nature versus being that perfection), \emph{self-subsisting Being} (whose entire essence is the act of existing), and \emph{diversified \dots\ more or less perfect} (being as an act realized in degrees, not a binary on/off toggle). Flatten ``participation'' into class membership and the causal step disappears.

\subsubsection*{Warrant}

The participation-causation rule is a \textbf{metaphysical principle}, illustrated from \textbf{experience} by iron ignited by fire. That God is self-subsisting Being and that subsisting being is one enter as \textbf{prior results} (q.\,3, a.\,4; q.\,11, aa.\,3--4, the second reinforced by a \textbf{metaphysical principle} about multiplication through recipients)---standing on \textbf{provenance only} until checked. The conclusion follows by \textbf{logical consequence}. The closing Plato and Aristotle citations supply \textbf{provenance only}: removing their names removes zero argumentative force.

\subsubsection*{Burden}

Granting q.\,3, a.\,4 and q.\,11, aa.\,3--4 makes ``all other beings exist by participation'' follow by exhaustion. The burden-bearing premise asserted on this page is the participation-causation principle: \textbf{whatever possesses existence non-essentially (by participation) must be efficiently caused by that to which existence belongs essentially.} An intelligent opponent can grant every definition here and still ask why possessing a feature non-essentially requires an efficient cause that \emph{is} that feature essentially---a wall is painted red by a painter who is not ``subsistent redness,'' and iron is heated by a stove that is not ``fire itself.''

\subsubsection*{Consequence}

Granting q.\,3, a.\,4, q.\,11, aa.\,3--4, and the participation-causation principle establishes that every being other than the one self-subsisting Being depends on that First Being for its existence. \emph{Not yet established:} (1) a temporal beginning of the world rather than eternal participated causation (q.\,46); (2) \emph{free} voluntary creation rather than necessary natural emanation like fire igniting iron (q.\,19 and q.\,46); and (3) the truth of the imported results in q.\,3, a.\,4 and q.\,11, aa.\,3--4 (tested in Exercise 7.4).

\subsubsection*{Denial}

\textbf{Rival framework} (not contradiction): treating existence as logical quantifier instantiation rather than a participated act in degrees, or accepting contingency without an essentially subsistent exemplar---leaving an \textbf{explanatory cost} if pressed to explain why contingent things exist at all.

\subsection{Exercises}

\subsubsection*{Exercise 7.1}
Here is Aquinas's Second Way (\emph{ST} I, q.\,2, a.\,3, corp., 1920 English Dominican translation), where the Latin speaks of \emph{causis efficientibus ordinatis} (``efficient causes following in order'') and invokes the rule \emph{remota autem causa, removetur effectus} (``to take away the cause is to take away the effect''):

\begin{quote}
The second way is from the nature of the efficient cause. In the world of sense we find there is an order of efficient causes. There is no case known (neither is it, indeed, possible) in which a thing is found to be the efficient cause of itself; for so it would be prior to itself, which is impossible. Now in efficient causes it is not possible to go on to infinity, because in all efficient causes following in order, the first is the cause of the intermediate cause, and the intermediate is the cause of the ultimate cause, whether the intermediate cause be several, or only one. Now to take away the cause is to take away the effect. Therefore, if there be no first cause among efficient causes, there will be no ultimate, nor any intermediate cause. But if in efficient causes it is possible to go on to infinity, there will be no first efficient cause, neither will there be an ultimate effect, nor any intermediate efficient causes; all of which is plainly false. Therefore it is necessary to admit a first efficient cause, to which everyone gives the name of God.
\end{quote}

Produce a full six-question analysis, answering explicitly where the ordered-series claim enters and whether this page itself tells you the series is ordered \emph{per se}.

\subsubsection*{Exercise 7.2}
Read this unannotated exposition of divine analogy:

\begin{quote}
Because God infinitely transcends creatures, no perfection-term like ``wisdom'' can be predicated univocally of God and creatures. Yet St.\ Thomas shows such words are not purely equivocal either, but analogical by causal proportion. Consequently, whenever reason shows that human wisdom implies truthful judgment, that term carries cleanly across a syllogism about God without risk of equivocation, because analogy guarantees a shared core meaning linking creature and Creator.
\end{quote}

Locate the exact defect and explain why it fails. Then repair the argument honestly \emph{without} smuggling the conclusion into a definition, testing your repair with an intelligent opponent and naming one tempting definitional cheat that fails.

\subsubsection*{Exercise 7.3}
Here is an unannotated passage from St.\ Thomas Aquinas, \emph{ST} I, q.\,14, a.\,1, corp., ``Whether There Is Knowledge in God?'' (1920 English Dominican translation):

\begin{quote}
In God there exists the most perfect knowledge. To prove this, we must note that intelligent beings are distinguished from non-intelligent beings in that the latter possess only their own form; whereas the intelligent being is naturally adapted to have also the form of some other thing; for the idea of the thing known is in the knower. Hence it is manifest that the nature of a non-intelligent being is more contracted and limited; whereas the nature of intelligent beings has a greater amplitude and extension; therefore the Philosopher says (De Anima iii) that ``the soul is in a sense all things.'' Now the contraction of the form comes from the matter. Hence, as we have said above (Q.\,7, A.\,1) forms according as they are the more immaterial, approach more nearly to a kind of infinity. Therefore it is clear that the immateriality of a thing is the reason why it is cognitive; and according to the mode of immateriality is the mode of knowledge. Hence it is said in \emph{De Anima} ii that plants do not know, because they are wholly material. But sense is cognitive because it can receive images free from matter, and the intellect is still further cognitive, because it is more separated from matter and unmixed, as said in \emph{De Anima} iii. Since therefore God is in the highest degree of immateriality as stated above (Q.\,7, A.\,1), it follows that He occupies the highest place in knowledge.
\end{quote}

Produce a complete six-question analysis unaided, using the One-Page Analysis Sheet (Appendix B).

\subsubsection*{Exercise 7.4}
Two sentences from the second argument below appeared in Lesson 1 under Claim and Meaning; now analyse the full three-argument \emph{corpus} of St.\ Thomas Aquinas, \emph{ST} I, q.\,3, a.\,4, ``Whether Essence and Existence Are the Same in God?'' (1920 English Dominican translation), unannotated:

\begin{quote}
God is not only His own essence, as shown in the preceding article, but also His own existence. This may be shown in several ways. First, whatever a thing has besides its essence must be caused either by the constituent principles of that essence (like a property that necessarily accompanies the species---as the faculty of laughing is proper to a man---and is caused by the constituent principles of the species), or by some exterior agent---as heat is caused in water by fire. Therefore, if the existence of a thing differs from its essence, this existence must be caused either by some exterior agent or by its essential principles. Now it is impossible for a thing's existence to be caused by its essential constituent principles, for nothing can be the sufficient cause of its own existence, if its existence is caused. Therefore that thing, whose existence differs from its essence, must have its existence caused by another. But this cannot be true of God; because we call God the first efficient cause. Therefore it is impossible that in God His existence should differ from His essence. Secondly, existence is that which makes every form or nature actual; for goodness and humanity are spoken of as actual, only because they are spoken of as existing. Therefore existence must be compared to essence, if the latter is a distinct reality, as actuality to potentiality. Therefore, since in God there is no potentiality, as shown above (A.\,1), it follows that in Him essence does not differ from existence. Therefore His essence is His existence. Thirdly, because, just as that which has fire, but is not itself fire, is on fire by participation; so that which has existence but is not existence, is a being by participation. But God is His own essence, as shown above (A.\,3); if, therefore, He is not His own existence He will be not essential, but participated being. He will not therefore be the first being---which is absurd. Therefore God is His own existence, and not merely His own essence.
\end{quote}

Produce a complete six-question analysis of the three arguments unaided, matching Appendix B.

\subsection{Solutions}

\subsubsection*{Solution to Exercise 7.1}
\emph{Claim and plain-language reading.} We must admit a first efficient cause (named God by everyone), since sensible things show an order of efficient causes, nothing causes itself, and infinite regress would remove intermediate and final effects. \emph{Meaning.} Everything turns on \emph{causis efficientibus ordinatis}, entering in observation and again in the no-regress step. Does the page specify \emph{per se} ordering? Record a \textbf{genuine indeterminacy}: on the first defensible reading (implicit \emph{per se}), \emph{remota autem causa, removetur effectus} applies only to concurrent derived causes, confirmed by \emph{ST} I, q.\,46, a.\,2, ad 7 and \emph{SCG} I, c.\,13, n.\,33; on the second (unstated on the page), q.\,2, a.\,3 never writes \emph{per se} or contrasts it with \emph{per accidens}. Your question settles which reading governs: Aquinas's broader doctrine favors the first; auditing this paragraph alone favors the second. \emph{Warrant.} \textbf{Experience} (observed order; absurdity of no effects), \textbf{metaphysical principles} (no self-causation; elimination rule, also a \textbf{prior result} in \emph{SCG} I, c.\,13, n.\,33), \textbf{logical consequence}, and nominal \textbf{definition} (``gives the name of God''). \emph{Burden.} That the sensible order of efficient causes is a derived-causing (\emph{per se}) hierarchy collapsing without a first member. \emph{Consequence.} Granted that burden, this ordered series ends in at least one uncaused efficient cause. \emph{Not yet established:} uniqueness across all series, and divine attributes (simplicity through unity in qq.\,3--11; divine knowledge in q.\,14). \emph{Denial.} \textbf{Rival framework} (successive intrinsically powered agents), not contradiction.

\subsubsection*{Solution to Exercise 7.2}
The defect sits in the final sentence: treating analogy---which denies one identical defining account---as though it automatically secured univocal syllogistic unity, hiding the four-term fallacy pressed by Scotus and later analyzed by Cajetan (\emph{De nominum analogia}, 1498). \emph{Honest repair:} conclude instead that ``real causal proportion licenses a \emph{quia} inference to that perfection's higher presence in God, while leaving open how far strict syllogistic deduction reaches across analogical terms without a minimal univocal core.'' An opponent can grant every definition here and still deny that causal proportion warrants valid inference without univocity, proving the repaired premise is substantive. A tempting definitional cheat defines an ``analogical term'' as ``a non-univocal term guaranteed to preserve syllogistic validity across God and creatures,'' smuggling the conclusion into vocabulary.

\subsubsection*{Solution to Exercise 7.3}
\emph{Claim and plain-language reading.} God has the most perfect knowledge, because knowing means receiving other forms, matter restricts a thing to its own form, immateriality is therefore the root of cognition, and God is supremely immaterial. \emph{Meaning.} Fix \emph{having another's form} (cognitive reception, not physical change), \emph{contraction by matter}, and \emph{immateriality} (freedom from material limit). \emph{Warrant.} Distinguishing knowers from non-knowers pairs \textbf{definition} and \textbf{experience}; \emph{De Anima} citations stand on \textbf{provenance only}; matter's contraction and God's supreme immateriality enter as \textbf{prior results} (q.\,7, a.\,1); that immateriality is \emph{the reason why} a thing knows is a bold \textbf{metaphysical principle}; the conclusion is a \textbf{logical consequence}. \emph{Burden.} Passing over q.\,7, a.\,1, the burden on this page is that \textbf{immateriality is of itself the sufficient cause making a being cognitive.} A hostile reader can grant God's immateriality yet ask why lacking matter automatically yields conscious knowledge rather than inert, non-conscious immaterial existence. \emph{Consequence.} Granted q.\,7, a.\,1 and that sufficiency bridge, God knows supremely. \emph{Not yet established:} knowledge of contingent futures, providence, or free will (later in q.\,14 and qq.\,19--22). \emph{Denial.} \textbf{Rival framework} (non-cognitive abstracta or naturalist functionalism), not contradiction.

\subsubsection*{Solution to Exercise 7.4}
\emph{Claim and plain-language reading.} Essence and existence are identical in God by three routes: (1) distinct existence would be caused by another, impossible for the first efficient cause; (2) distinct essence stands to existence as potency to act, impossible without potentiality; (3) having existence without being existence makes a thing participated rather than the first being. \emph{Meaning.} Fix \emph{essence}, \emph{existence} (\emph{esse} as the real act actualizing every form, not mere propositional truth, which Reply 2 sets apart), \emph{act/potency}, and \emph{participation}. \emph{Warrant.} Each route joins a \textbf{metaphysical principle} (illustrated from \textbf{experience} by laughter, heated water, and fire) with a \textbf{prior result} (q.\,2, a.\,3; q.\,3, a.\,1; q.\,3, a.\,3---each \textbf{provenance only} until audited), concluding by \textbf{logical consequence}. \emph{Burden.} Underlying all three routes is the shared commitment that \textbf{existence (\emph{esse}) is a real metaphysical act entering into genuine composition with essence in finite things, rather than merely the logical fact that a concept is instantiated.} An opponent can grant every definition and deny that existence is a component act. \emph{Consequence.} Granted that ontology and the three prior results, essence and existence coincide in the First Cause. \emph{Not yet established:} (1) comprehension of \emph{what} God's act of existence is (explicitly denied in Reply 2); (2) uniqueness (q.\,11, aa.\,3--4); or (3) intelligence, freedom, or creation (qq.\,14, 19, 44--46). \emph{Denial.} \textbf{Rival framework} (predicate-logic quantification or scholastic critiques of real essence-existence composition), not contradiction.
